\section{Applications Beyond Pure Mathematics}

Knot theory appears whenever long flexible objects become entangled. Examples include:

\begin{itemize}[label=$\bigoplus$]
    \item circular DNA molecules and the action of topoisomerase enzymes;
    \item polymer chains and molecular entanglement;
    \item vortex lines in fluids and plasmas;
    \item knot-based models in quantum computation
\end{itemize}

For two oriented components C1 and C2 in a link diagram, the linking number can be calculated from crossings between the two components:
\begin{equation}
    \text{lk}(C_1, C_2) = \frac{1}{2} \sum_{p\in C_1 \cap C_2} \varepsilon(P).
\end{equation}

This quantity measures how many times the two components wind around one another, with orientation taken
into account.

For a link $L= C_1 \cap C_2$, the pairwise linking numbers may be collected into the symmetric linking matrix

\begin{equation}
    \Lambda (L) = [\lambda_{ij}]_{i,j=1}^n , \quad \lambda_{ij} = \begin{cases}
      \text{lk}(C_1, C_2), & i\neq j, \\
      0, & i=j
    \end{cases}
\end{equation}
